% Bagian: first-order-logic
% Bab: axiomatic-deduction
% Subbagian: provability-consistency

\documentclass[../../../include/open-logic-section]{subfiles}

\begin{document}

\iftag{FOL}
      {\olfileid[id]{fol}{axd}{prv}}
      {\olfileid[id]{pl}{axd}{prv}}

\olsection{\usetoken{S}{derivability} dan Konsistensi}

Sekarang kita akan membuktikan sejumlah sifat relasi !!{derivability}.
Sifat-sifat tersebut menarik secara mandiri, tetapi masing-masing akan
berperan dalam pembuktian teorema kelengkapan.

\begin{prop}\ollabel{prop:provability-contr}
  Jika $\Gamma \Proves !A$ dan $\Gamma \cup \{!A\}$ tak konsisten,
  maka $\Gamma$ tak konsisten.
\end{prop}

\begin{proof}
  Jika $\Gamma \cup \{!A\}$ tak konsisten, maka
  $\Gamma \cup \{!A\} \Proves \lfalse$. Menurut
  \olref[ptn]{prop:reflexivity}, $\Gamma \Proves !B$ untuk setiap
  $!B \in \Gamma$. Karena menurut hipotesis kita juga mempunyai
  $\Gamma \Proves !A$, maka $\Gamma \Proves !B$ untuk setiap
  $!B \in \Gamma \cup \{!A\}$. Menurut
  \olref[ptn]{prop:transitivity}, $\Gamma \Proves \lfalse$; yakni,
  $\Gamma$ tak konsisten.
\end{proof}

\begin{prop}
\ollabel{prop:prov-incons}
$\Gamma \Proves !A$ jika dan hanya jika $\Gamma \cup \{\lnot !A\}$
tak konsisten.
\end{prop}

\begin{proof}
Pertama, andaikan $\Gamma \Proves !A$. Maka menurut
\olref[ptn]{prop:monotonicity},
$\Gamma \cup \{\lnot !A\} \Proves !A$. Menurut
\olref[ptn]{prop:reflexivity},
$\Gamma \cup \{\lnot !A\} \Proves \lnot !A$. Kita juga mempunyai
$\Proves \lnot !A \lif (!A \lif \lfalse)$ menurut
\olref[prp]{ax:lnot2}. Jadi, melalui dua penerapan
\olref[ded]{prop:mp}, kita memperoleh
$\Gamma \cup \{\lnot !A\} \Proves \lfalse$.

Sekarang andaikan $\Gamma \cup \{\lnot !A\}$ tak konsisten; yakni,
$\Gamma \cup \{\lnot !A\} \Proves \lfalse$. Menurut teorema deduksi,
$\Gamma \Proves \lnot !A \lif \lfalse$. Menurut
\olref[prp]{ax:lfalse2},
$\Gamma \Proves (\lnot !A \lif \lfalse) \lif \lnot\lnot !A$,
sehingga \olref[ded]{prop:mp} memberikan
$\Gamma \Proves \lnot\lnot !A$. Karena
$\Gamma \Proves \lnot\lnot !A \lif !A$
(\olref[prp]{ax:dne}), penerapan \olref[ded]{prop:mp} sekali lagi
memberikan $\Gamma \Proves !A$.
\end{proof}

\begin{prob}
Buktikan bahwa $\Gamma \Proves \lnot !A$ jika dan hanya jika
$\Gamma \cup \{!A\}$ tak konsisten.
\end{prob}

\begin{prop}\ollabel{prop:explicit-inc}
  Jika $\Gamma \Proves !A$ dan $\lnot !A \in \Gamma$, maka $\Gamma$
  tak konsisten.
\end{prop}

\begin{proof}
  Menurut \olref[prp]{ax:lnot2},
  $\Gamma \Proves \lnot !A \lif (!A \lif \lfalse)$.
  Dua penerapan \olref[ded]{prop:mp} memberikan
  $\Gamma \Proves \lfalse$.
\end{proof}

\begin{prop}\ollabel{prop:provability-exhaustive}
  Jika $\Gamma \cup \{!A\}$ dan $\Gamma \cup \{\lnot !A\}$ keduanya
  tak konsisten, maka $\Gamma$ tak konsisten.
\end{prop}

\begin{proof}
Latihan.
\end{proof}

\tagprob{FOL}
\begin{prob}
  Buktikan \olref[fol][axd][prv]{prop:provability-exhaustive}.
\end{prob}
\tagendprob

\tagprob{notFOL}
\begin{prob}
  Buktikan \olref[pl][axd][prv]{prop:provability-exhaustive}.
\end{prob}
\tagendprob

\end{document}
